\chapter{Proof Techniques}

A proof is a finite chain of statements, each of which is an assumption, a previously proved
result, or follows logically from earlier lines, ending in the statement to be proved. This
chapter collects the standard proof templates used throughout the course.

\section{Proving universal and existential statements}

\begin{itemize}
  \item To prove $\forall x \in S,\ P(x)$: let $x$ be an \emph{arbitrary} element of $S$ (do not
        pick a specific one) and show $P(x)$.
  \item To prove $\exists x \in S,\ P(x)$: exhibit (construct) a specific $x \in S$ and verify
        $P(x)$. You do not need to explain how you found it.
  \item To \emph{disprove} $\forall x \in S,\ P(x)$: give a single \emph{counterexample}
        $x \in S$ with $P(x)$ false.
  \item To disprove $\exists x \in S,\ P(x)$: prove $\forall x \in S,\ \neg P(x)$.
\end{itemize}

\begin{example}
Prove there is an integer $n$ such that $n^2 - 5n + 6 = 0$.
\emph{Proof.} Take $n = 2$: $4 - 10 + 6 = 0$. \qed
\end{example}

\begin{example}[Counterexample]
Disprove: ``for all real $a,b$, if $a^2 = b^2$ then $a = b$''. Take $a = 1$, $b = -1$: then
$a^2 = b^2 = 1$ but $a \neq b$.
\end{example}

\section{Direct proof}

To prove $A \Rightarrow B$ directly, assume $A$ and deduce $B$. Many proofs in this course
start from definitions, so it pays to state them exactly.

\begin{definition}
An integer $n$ is \emph{even} if $n = 2k$ for some $k \in \Z$, and \emph{odd} if $n = 2k+1$ for
some $k\in\Z$. For integers $a,b$, we say $a$ \emph{divides} $b$, written $a \mid b$, if there
is $k \in \Z$ with $b = ka$.
\end{definition}

\begin{proposition}
If $n$ is an odd integer then $n^2$ is odd.
\end{proposition}
\begin{proof}
Let $n$ be odd, so $n = 2k + 1$ for some $k \in \Z$. Then
$n^2 = 4k^2 + 4k + 1 = 2(2k^2 + 2k) + 1$. Since $2k^2+2k \in \Z$, $n^2$ is odd.
\end{proof}

\begin{proposition}
For all $a,b,c \in \Z$, if $a \mid b$ and $a \mid c$ then $a \mid (b + c)$.
\end{proposition}
\begin{proof}
Since $a \mid b$ and $a \mid c$ there are integers $k,\ell$ with $b = ka$ and $c = \ell a$.
Then $b + c = (k + \ell)a$ with $k+\ell \in \Z$, so $a \mid (b+c)$.
\end{proof}

\begin{important}[Working backwards is not a proof]
When solving, it is fine to start from the conclusion and work backwards on scratch paper,
but the written proof must go from hypothesis to conclusion with each step justified.
In particular, ``$B \Rightarrow \TT$'' does not prove $B$: from $-1 = 1$ one gets $1 = 1$ by
squaring.
\end{important}

\begin{example}[AM--GM]
Prove that for all real $x,y \ge 0$, $\sqrt{xy} \le \frac{x+y}{2}$.

\emph{Scratch work.} $\sqrt{xy} \le \frac{x+y}{2} \iff 2\sqrt{xy} \le x + y \iff 0 \le x - 2\sqrt{xy}
+ y = (\sqrt{x} - \sqrt{y})^2$.

\emph{Proof.} Since $\sqrt x, \sqrt y$ are real, $(\sqrt x - \sqrt y)^2 \ge 0$, i.e.\
$x - 2\sqrt{xy} + y \ge 0$. Rearranging, $\sqrt{xy} \le \frac{x+y}{2}$. \qed
\end{example}

\section{Proof by cases}

If $A_1 \vee A_2 \vee \cdots \vee A_k$ is known to hold (the cases are exhaustive) and each
$A_i \Rightarrow B$, then $B$ holds.

\begin{proposition}
For every integer $n$, $n^2 + n$ is even.
\end{proposition}
\begin{proof}
\emph{Case 1: $n$ even}, $n = 2k$. Then $n^2 + n = 4k^2 + 2k = 2(2k^2 + k)$, even.
\emph{Case 2: $n$ odd}, $n = 2k+1$. Then $n^2 + n = (2k+1)(2k+2) = 2(2k+1)(k+1)$, even.
Every integer is even or odd, so the result holds.
\end{proof}

\section{Proof by contrapositive}

Since $A \Rightarrow B \equiv \neg B \Rightarrow \neg A$, we may instead assume $\neg B$ and
derive $\neg A$. This helps when $\neg B$ gives something concrete to work with.

\begin{proposition}
Let $n \in \Z$. If $n^2$ is even, then $n$ is even.
\end{proposition}
\begin{proof}
We prove the contrapositive: if $n$ is odd then $n^2$ is odd. This was proved above.
\end{proof}

\begin{proposition}
For $x \in \R$, if $x^3 + 5x < 6$ then $x < 1$.
\end{proposition}
\begin{proof}
Contrapositive: suppose $x \geq 1$. Then $x^3 \ge 1$ and $5x \ge 5$, so $x^3 + 5x \ge 6$.
\end{proof}

\section{Proof by contradiction}

To prove $B$, assume $\neg B$ and derive a contradiction (a statement of the form $C\wedge
\neg C$). Then $\neg B$ must be false, so $B$ is true.

\begin{theorem}
$\sqrt 2$ is irrational.
\end{theorem}
\begin{proof}
Suppose, for a contradiction, that $\sqrt 2 = p/q$ with $p,q \in \Z$, $q > 0$, and the fraction
in lowest terms (so $p$ and $q$ are not both even). Then $p^2 = 2q^2$, so $p^2$ is even, hence
$p$ is even by the proposition above. Write $p = 2k$; then $4k^2 = 2q^2$, so $q^2 = 2k^2$ is
even and $q$ is even. Both $p$ and $q$ are even, a contradiction.
\end{proof}

\begin{theorem}[Euclid]
There are infinitely many primes.
\end{theorem}
\begin{proof}
Suppose there are finitely many, $p_1,\dots,p_k$. Let $N = p_1p_2\cdots p_k + 1 > 1$. $N$ has
some prime divisor (every integer $>1$ does, Chapter 5), which must be some $p_i$. Then $p_i
\mid N$ and $p_i \mid p_1 \cdots p_k$, so $p_i \mid N - p_1\cdots p_k = 1$, which is impossible.
\end{proof}

\begin{aside}
Contradiction versus contrapositive: to prove $A \Rightarrow B$ by contradiction we assume $A$
and $\neg B$ and reach \emph{any} contradiction; the contrapositive is the special case where
we derive $\neg A$.
\end{aside}

\section{If and only if}

To prove $A \Leftrightarrow B$ prove both directions, $A \Rightarrow B$ and $B \Rightarrow A$
(each by any method). Alternatively build a chain of equivalences
$A \iff C_1 \iff \cdots \iff B$, being careful that every step really is reversible.

\begin{proposition}
Let $n \in \Z$. Then $n$ is even if and only if $n^2$ is even.
\end{proposition}
\begin{proof}
($\Rightarrow$) If $n = 2k$ then $n^2 = 2(2k^2)$ is even.
($\Leftarrow$) Proved above by contrapositive.
\end{proof}

\begin{example}
Prove that for $x\in\R$, $x^2 - 3x + 2 < 0 \iff 1 < x < 2$.
We have $x^2 - 3x + 2 = (x-1)(x-2)$. A product of two reals is negative if and only if the
factors have opposite signs. $x - 1 < 0 < x - 2$ is impossible (it would need $x < 1$ and
$x > 2$), so the product is negative iff $x - 1 > 0 > x - 2$, i.e.\ iff $1 < x < 2$.
\end{example}

\section{Uniqueness}

To prove ``there is a unique $x \in S$ with $P(x)$'' (written $\exists!\, x$):
\begin{enumerate}
  \item \emph{Existence}: find some $x$ with $P(x)$.
  \item \emph{Uniqueness}: assume $P(x)$ and $P(y)$ and show $x = y$; or assume $x \ne y$
        both satisfy $P$ and derive a contradiction.
\end{enumerate}

\begin{proposition}
For every $a \in \R$ and $b \neq 0$, the equation $bx = a$ has a unique real solution.
\end{proposition}
\begin{proof}
\emph{Existence:} $x = a/b$ works since $b(a/b) = a$.
\emph{Uniqueness:} if $bx = a$ and $by = a$ then $b(x - y) = 0$; since $b \ne 0$, $x = y$.
\end{proof}

\section{Common pitfalls}
\begin{itemize}
  \item Proving a universal statement by checking examples. Examples can only disprove.
  \item Using the same letter for two different unknowns: if $m$ and $n$ are both odd, write
        $m = 2k + 1$ and $n = 2\ell + 1$, not both with $k$.
  \item Dividing by something that might be zero.
  \item In a contradiction proof, forgetting to say what the contradiction is.
\end{itemize}

\section{Summary}
\begin{itemize}
  \item $\forall$: take an arbitrary element. $\exists$: construct a witness. To disprove
        $\forall$: one counterexample.
  \item Direct proof: assume the hypothesis, reach the conclusion, using definitions.
  \item Contrapositive: prove $\neg B \Rightarrow \neg A$. Contradiction: assume the
        statement is false and derive an impossibility.
  \item $A \Leftrightarrow B$: prove both directions. Uniqueness: existence plus ``any two are
        equal''.
\end{itemize}
